\section{Main Results}

We organize the empirical findings around six research questions.

\paragraph{RQ1: Can CLI agents complete end-to-end research under minimal scaffolding?}
Yes, in the limited sense of producing complete research-like artifacts. All three agents generate papers, code, logs, results, and reviews across the full domain set, yielding 117 papers. However, completion should not be confused with validity. Many papers contain incomplete baselines, unsupported claims, setting mismatches, or results that cannot be reconstructed from artifacts.

\begin{table}[t]
\centering
\small
\begin{tabular}{lrrrrrr}
\toprule
System & $n$ & \sar\ mean & \sar\ std. & Min & Max & $\geq 5.0$ \\
\midrule
Claude Code & 39 & 5.45 & 0.70 & 3.1 & 6.3 & 82\% \\
FARS & 102 & 5.06 & 0.62 & 3.0 & 6.3 & 76\% \\
Codex & 39 & 4.93 & 0.85 & 2.7 & 6.3 & 64\% \\
Kimi Code & 39 & 4.24 & 0.84 & 2.0 & 5.3 & 26\% \\
\bottomrule
\end{tabular}
\caption{PDF-only \sar\ comparison. We treat FARS as a secondary external reference point rather than as the central comparison target.}
\label{tab:sar-system}
\end{table}

\begin{figure}[t]
\centering
\includegraphics[width=0.86\linewidth]{sar_score_distributions.png}
\caption{\sar\ score distributions. PDF-only review assigns many generated papers scores near the middle of the ICLR scale, with Claude Code highest among the three CLI agents.}
\label{fig:sar-dist}
\end{figure}

\paragraph{RQ2: How do \sar\ and \pr\ differ?}
Table~\ref{tab:sar-system} and Figure~\ref{fig:sar-dist} show the PDF-only \sar\ view. Claude Code scores 5.45, Codex 4.93, and Kimi Code 4.24; FARS reports 5.06 on 102 papers. A human-calibration sample gives ICLR accepted papers a \sar\ mean of 5.59, a weighted ICLR mean of 5.42, and rejected papers 5.34. Human-annotated \sar\ accept rates are more separated: ICLR accepted 76.0\%, ICLR rejected 52.0\%, Claude Code 41.0\%, FARS 21.6\%, Codex 12.8\%, and Kimi Code 5.1\%.

Artifact-grounded \pr\ is stricter: Claude Code scores 4.59, Codex 4.51, and Kimi Code 3.38. The important point is not that \pr\ is a better scalar judge in all settings, but that it evaluates a different object. \sar\ evaluates the final PDF. \pr\ evaluates whether the PDF is supported by the research artifacts.

\begin{figure}[t]
\centering
\begin{subfigure}{0.48\linewidth}
\centering
\includegraphics[width=\linewidth]{pr_score_distributions.png}
\caption{\pr\ distributions}
\label{fig:pr-dist}
\end{subfigure}
\hfill
\begin{subfigure}{0.48\linewidth}
\centering
\includegraphics[width=\linewidth]{internal_vs_stanford_scatter.png}
\caption{\sar\ versus \pr}
\label{fig:sar-pr-scatter}
\end{subfigure}
\caption{Artifact-grounded review shifts the evaluation target from paper plausibility to paper-artifact faithfulness. Scores fall most when artifacts reveal unsupported results, setting mismatches, or incomplete experiments.}
\label{fig:pr-and-scatter}
\end{figure}

\paragraph{RQ3: What are the dominant bottlenecks?}
Across agents, experimental rigor and significance are the recurring weaknesses. Figure~\ref{fig:dimensions} shows that Codex leads on reproducibility, reference integrity, and results integrity, while Claude Code leads on novelty and significance. Kimi Code trails on most dimensions. The lowest dimensions are not simply novelty: the agents frequently produce plausible ideas but fail to validate them with strong baselines, adequate scale, consistent protocols, or faithful reporting.

\begin{figure}[t]
\centering
\includegraphics[width=0.9\linewidth]{comparison_3way_dimensions.png}
\caption{Per-dimension \pr\ scores. Codex is strongest on integrity-oriented dimensions, Claude Code is strongest on novelty and significance, and experimental rigor remains the dominant bottleneck.}
\label{fig:dimensions}
\end{figure}

\begin{table}[t]
\centering
\small
\begin{tabular}{lrrr}
\toprule
Integrity category & Claude Code & Codex & Kimi Code \\
\midrule
Results mismatch only & 6/39 & 2/39 & 4/39 \\
Setting mismatch only & 10/39 & 1/39 & 5/39 \\
Both results and setting mismatch & 12/39 & 2/39 & 30/39 \\
Fake reference & 14/39 & 3/39 & 28/39 \\
\bottomrule
\end{tabular}
\caption{Manual integrity audit. Artifact inspection reveals paper-artifact mismatches and fake references that are often not visible from the PDF alone.}
\label{tab:integrity}
\end{table}

\paragraph{RQ4: Do agents exhibit distinct research personas?}
The agents show stable differences in research style (Table~\ref{tab:personas}). Claude Code behaves like a full-stack researcher: it writes the longest papers, uses more figures and tables, and balances empirical studies with novel methods. Codex behaves like an empirical scientist: 87\% of its papers are controlled empirical studies, with high integrity but lower novelty. Kimi Code behaves like a system builder: it proposes acronym-heavy named methods, but has high mismatch and fake-reference rates.

\begin{table}[t]
\centering
\small
\begin{tabularx}{\linewidth}{lrrrrX}
\toprule
Agent & Novel method & Benchmark & Extension & Empirical & Summary \\
\midrule
Claude Code & 33\% & 8\% & 13\% & 46\% & Balanced, broad, longest papers, more figures and tables \\
Codex & 8\% & 0\% & 5\% & 87\% & Controlled empirical studies, highest integrity, lower novelty \\
Kimi Code & 79\% & 10\% & 0\% & 10\% & Named methods and acronyms, high mismatch and fake-reference rates \\
\bottomrule
\end{tabularx}
\caption{Research type distribution and qualitative personas. Style explains part of the score profile: Codex is reliable but narrow, Claude Code is broader, and Kimi Code often overproduces method framing relative to evidence.}
\label{tab:personas}
\end{table}

\paragraph{RQ5: Does more compute improve research quality?}
More compute does not reliably improve quality in this study. Figure~\ref{fig:domains-time} shows CPU/GPU differences and per-domain variation. Under \pr, all agents score higher on CPU tasks than GPU tasks, while \sar\ sometimes shows the opposite trend. This suggests that GPU domains can yield more polished-looking papers while also being harder to execute faithfully.

\begin{figure}[t]
\centering
\begin{subfigure}{0.48\linewidth}
\centering
\includegraphics[width=\linewidth]{pr_per_domain_breakdown.png}
\caption{Per-domain \pr}
\end{subfigure}
\hfill
\begin{subfigure}{0.48\linewidth}
\centering
\includegraphics[width=\linewidth]{comparison_3way_cpu_gpu_gap.png}
\caption{CPU/GPU gap}
\end{subfigure}
\caption{Domain and hardware patterns. Artifact-grounded review exposes a CPU/GPU divergence: visually polished GPU-domain papers often suffer from harder execution and verification conditions.}
\label{fig:domains-time}
\end{figure}

\begin{table}[t]
\centering
\small
\begin{tabular}{lrrr}
\toprule
GPU domain seed & Codex H100 \pr & Codex A6000 \pr & Difference \\
\midrule
AI for Biology & 4.89 & 4.67 & +0.22 \\
Interpretability of Learned Repr. & 4.78 & 4.67 & +0.11 \\
Computer Vision & 4.33 & 4.33 & +0.00 \\
Natural Language Processing & 4.22 & 4.33 & -0.11 \\
Privacy in ML & 4.22 & 4.00 & +0.22 \\
Supervised Repr. Learning & 4.22 & 4.67 & -0.45 \\
Generative Models & 3.78 & 4.33 & -0.55 \\
Datasets and Benchmarks & 3.67 & 4.67 & -1.00 \\
\midrule
Average & 4.26 & 4.50 & -0.24 \\
\bottomrule
\end{tabular}
\caption{H100 scaling experiment for Codex GPU runs. Moving from A6000 to H100 does not improve aggregate \pr; the average change is about -0.25 points. This suggests, but does not prove, that compute is not the primary bottleneck.}
\label{tab:h100}
\end{table}

\paragraph{RQ6: What failures are revealed by artifact inspection?}
Table~\ref{tab:integrity} summarizes manual integrity failures. Kimi Code has the highest mismatch rates: 30/39 papers exhibit both result and setting mismatch, and 28/39 contain fake references. Claude Code is intermediate, while Codex has the lowest mismatch and fake-reference rates. These defects are precisely the type that can survive PDF-only review. A paper can present a consistent story while its artifacts show a smaller benchmark sweep, a different evaluation protocol, a missing baseline, or numbers copied from the wrong file.

\begin{figure}[t]
\centering
\begin{subfigure}{0.48\linewidth}
\centering
\includegraphics[width=\linewidth]{comparison_3way_reviewer_scores.png}
\caption{Reviewer severity}
\end{subfigure}
\hfill
\begin{subfigure}{0.48\linewidth}
\centering
\includegraphics[width=\linewidth]{comparison_3way_time_per_stage.png}
\caption{Time by stage}
\end{subfigure}
\caption{Review and runtime patterns. Codex is the strictest reviewer, Kimi Code the most lenient, motivating multi-reviewer aggregation. Experiments dominate wall time, so failures are not simply due to agents skipping the experimental stage.}
\label{fig:reviewer-time}
\end{figure}

\begin{figure}[t]
\centering
\includegraphics[width=0.72\linewidth]{comparison_3way_wall_time.png}
\caption{Wall-time distribution. Claude Code takes longest on average, Kimi Code is fastest, and the long tail is driven by experiment and revision loops.}
\label{fig:walltime}
\end{figure}
